\documentclass[a4paper]{article}
\usepackage[T1]{fontenc}
\usepackage{lmodern}
\usepackage[margin=18mm]{geometry}
\usepackage[
  columns={Backlog,Doing,Review,Done},
  columnsep=2mm,
  reference-date=2026-07-18
]{kanban}
\pagestyle{empty}
\kbdefinepalette{release}{kanbanBlue,kanbanTeal,kanbanAmber,kanbanPlum}
\begin{document}
\section*{Release board}
% The reference date makes deadline classification reproducible. Calendar
% deadlines use the default US form MM/DD/YYYY.
\begin{kanban}[palette=release,mode=light]
  % Tasks may be declared in any cross-column order. Within each column they
  % retain their source order.
  \kanbantask[prio=high,author=Alice,deadline=01/01/2000]
    {Backlog}{Finalize the public API}
  \kanbantask[prio=medium,author=Bob,
    due=07/18/2026]
    {Doing}{Write regression tests}
  \kanbantask[prio=low,author=Carol,deadline=12/31/2099]
    {Review}{Proofread the manual}
  % \kbtask is an exact concise alias when that name was free at package load.
  \kbtask[author=Alice,completed]{Done}{Create the repository}
\end{kanban}
\bigskip
\section*{Dark board}
% The longest heading determines the height of every heading cell. The empty
% column still extends to the same bottom edge as the populated columns.
\begin{kanban}[
  columns={Planned,Active review,Done,Empty},
  palette=muted,
  mode=dark,
  columnsep=3mm
]
  \kanbantask[prio=high,author=Alice,deadline=01/01/2000]
    {Planned}{Prepare the release archive}
  \kanbantask[prio=medium,author=Bob,
    due=07/18/2026]
    {Active review}{Run the regression suite}
  \kanbantask[completed,deadline=12/31/2099]
    {Done}{Document the palette API}
\end{kanban}
\end{document}
